\PassOptionsToPackage{unicode}{hyperref}
\PassOptionsToPackage{hyphens}{url}
\PassOptionsToPackage{dvipsnames,svgnames,x11names}{xcolor}
\documentclass[10pt,letterpaper]{article}

\usepackage{xcolor}
\usepackage[margin=2.0cm]{geometry}
\usepackage{lmodern}
\usepackage{amsmath,amssymb,amsfonts,mathtools}
\usepackage{textcomp}
\usepackage{microtype}
\usepackage{parskip}
\usepackage{longtable,booktabs,array,calc}
\usepackage{caption}
\captionsetup[table]{skip=6pt,font=small,labelfont=bf}
\captionsetup[figure]{font=small,labelfont=bf}
\newcounter{none}
\usepackage{etoolbox}
\usepackage{graphicx}
\usepackage{fancyvrb}
\newcommand{\VerbBar}{|}
\newcommand{\VERB}{\Verb[commandchars=\\\{\}]}
\DefineVerbatimEnvironment{Highlighting}{Verbatim}{commandchars=\\\{\},fontsize=\footnotesize}
\RecustomVerbatimEnvironment{verbatim}{Verbatim}{fontsize=\footnotesize}
\newenvironment{Shaded}{}{}
\newcommand{\AlertTok}[1]{\textcolor[rgb]{0.80,0.00,0.00}{\textbf{#1}}}
\newcommand{\AnnotationTok}[1]{\textcolor[rgb]{0.38,0.63,0.69}{\textbf{\textit{#1}}}}
\newcommand{\AttributeTok}[1]{\textcolor[rgb]{0.49,0.56,0.16}{#1}}
\newcommand{\BaseNTok}[1]{\textcolor[rgb]{0.25,0.63,0.44}{#1}}
\newcommand{\BuiltInTok}[1]{\textcolor[rgb]{0.00,0.50,0.00}{#1}}
\newcommand{\CharTok}[1]{\textcolor[rgb]{0.25,0.44,0.63}{#1}}
\newcommand{\CommentTok}[1]{\textcolor[rgb]{0.38,0.63,0.69}{\textit{#1}}}
\newcommand{\CommentVarTok}[1]{\textcolor[rgb]{0.38,0.63,0.69}{\textbf{\textit{#1}}}}
\newcommand{\ConstantTok}[1]{\textcolor[rgb]{0.53,0.00,0.00}{#1}}
\newcommand{\ControlFlowTok}[1]{\textcolor[rgb]{0.00,0.44,0.13}{\textbf{#1}}}
\newcommand{\DataTypeTok}[1]{\textcolor[rgb]{0.56,0.13,0.00}{#1}}
\newcommand{\DecValTok}[1]{\textcolor[rgb]{0.25,0.63,0.44}{#1}}
\newcommand{\DocumentationTok}[1]{\textcolor[rgb]{0.73,0.13,0.13}{\textit{#1}}}
\newcommand{\ErrorTok}[1]{\textcolor[rgb]{0.80,0.00,0.00}{\textbf{#1}}}
\newcommand{\ExtensionTok}[1]{#1}
\newcommand{\FloatTok}[1]{\textcolor[rgb]{0.25,0.63,0.44}{#1}}
\newcommand{\FunctionTok}[1]{\textcolor[rgb]{0.02,0.16,0.49}{#1}}
\newcommand{\ImportTok}[1]{\textcolor[rgb]{0.00,0.50,0.00}{\textbf{#1}}}
\newcommand{\InformationTok}[1]{\textcolor[rgb]{0.38,0.63,0.69}{\textbf{\textit{#1}}}}
\newcommand{\KeywordTok}[1]{\textcolor[rgb]{0.00,0.44,0.13}{\textbf{#1}}}
\newcommand{\NormalTok}[1]{#1}
\newcommand{\OperatorTok}[1]{\textcolor[rgb]{0.40,0.40,0.40}{#1}}
\newcommand{\OtherTok}[1]{\textcolor[rgb]{0.00,0.44,0.13}{#1}}
\newcommand{\PreprocessorTok}[1]{\textcolor[rgb]{0.74,0.48,0.00}{#1}}
\newcommand{\RegionMarkerTok}[1]{#1}
\newcommand{\SpecialCharTok}[1]{\textcolor[rgb]{0.25,0.44,0.63}{#1}}
\newcommand{\SpecialStringTok}[1]{\textcolor[rgb]{0.73,0.40,0.53}{#1}}
\newcommand{\StringTok}[1]{\textcolor[rgb]{0.25,0.44,0.63}{#1}}
\newcommand{\VariableTok}[1]{\textcolor[rgb]{0.10,0.09,0.49}{#1}}
\newcommand{\VerbatimStringTok}[1]{\textcolor[rgb]{0.25,0.44,0.63}{#1}}
\newcommand{\WarningTok}[1]{\textcolor[rgb]{0.38,0.63,0.69}{\textbf{\textit{#1}}}}

\usepackage{newunicodechar}
\newunicodechar{µ}{\ifmmode\mu\else\textmu\fi}
\newunicodechar{°}{\ifmmode^\circ\else\textdegree\fi}
\newunicodechar{·}{\ifmmode\cdot\else\textperiodcentered\fi}
\newunicodechar{²}{\ensuremath{^{2}}}
\newunicodechar{³}{\ensuremath{^{3}}}
\newunicodechar{¹}{\ensuremath{^{1}}}
\newunicodechar{⁴}{\ensuremath{^{4}}}
\newunicodechar{⁶}{\ensuremath{^{6}}}
\newunicodechar{⁻}{\ensuremath{^{-}}}
\newunicodechar{₀}{\ensuremath{_{0}}}
\newunicodechar{ₙ}{\ensuremath{_{n}}}
\newunicodechar{Δ}{\ensuremath{\Delta}}
\newunicodechar{Ω}{\ensuremath{\Omega}}
\newunicodechar{α}{\ensuremath{\alpha}}
\newunicodechar{β}{\ensuremath{\beta}}
\newunicodechar{δ}{\ensuremath{\delta}}
\newunicodechar{ε}{\ensuremath{\varepsilon}}
\newunicodechar{θ}{\ensuremath{\theta}}
\newunicodechar{κ}{\ensuremath{\kappa}}
\newunicodechar{λ}{\ensuremath{\lambda}}
\newunicodechar{π}{\ensuremath{\pi}}
\newunicodechar{ρ}{\ensuremath{\rho}}
\newunicodechar{σ}{\ensuremath{\sigma}}
\newunicodechar{τ}{\ensuremath{\tau}}
\newunicodechar{→}{\ensuremath{\rightarrow}}
\newunicodechar{∎}{\ensuremath{\blacksquare}}
\newunicodechar{−}{\ensuremath{-}}
\newunicodechar{∧}{\ensuremath{\wedge}}
\newunicodechar{∫}{\ensuremath{\int}}
\newunicodechar{≈}{\ensuremath{\approx}}
\newunicodechar{≤}{\ensuremath{\leq}}
\newunicodechar{≥}{\ensuremath{\geq}}
\newunicodechar{⊆}{\ensuremath{\subseteq}}
\newunicodechar{✓}{\ensuremath{\checkmark}}
\newunicodechar{≠}{\ensuremath{\neq}}
\newunicodechar{↔}{\ensuremath{\leftrightarrow}}
\newunicodechar{≪}{\ensuremath{\ll}}
\newunicodechar{↛}{\ensuremath{\nrightarrow}}
\newunicodechar{⇒}{\ensuremath{\Rightarrow}}
\newunicodechar{∪}{\ensuremath{\cup}}
\newunicodechar{⟨}{\ensuremath{\langle}}
\newunicodechar{⟩}{\ensuremath{\rangle}}
\newunicodechar{・}{\textperiodcentered}

\usepackage{xurl}
\usepackage{bookmark}
\usepackage{hyperref}

\setcounter{secnumdepth}{-\maxdimen}
\providecommand{\tightlist}{\setlength{\itemsep}{1pt}\setlength{\parskip}{0pt}}

\definecolor{openieAmber}{HTML}{D97706}
\definecolor{openieAmberBg}{HTML}{FFFBEB}
\definecolor{openieAmberBorder}{HTML}{F59E0B}
\definecolor{openieSlate}{HTML}{334155}
\definecolor{openieDark}{HTML}{0F172A}
\definecolor{openieCodeBg}{HTML}{F8FAFC}
\definecolor{openieCodeBorder}{HTML}{CBD5E1}
\definecolor{openieBlue}{HTML}{0057B8}

\hypersetup{
  pdftitle={Metabolic Intelligence: Budget-Native Compute from Tag to Campus},
  pdfauthor={David Charlot, Open Interface Engineering},
  colorlinks=true,
  linkcolor=openieBlue,
  citecolor=openieBlue,
  urlcolor=openieBlue,
  pdfcreator={XeTeX / Tectonic — Open Interface Engineering}
}

\urlstyle{same}
\setlength{\emergencystretch}{3em}
\hfuzz=6pt
\sloppy

\newcommand{\openiefig}[3]{%
  \IfFileExists{#1}{%
    \includegraphics[width=#2\linewidth,height=0.70\textheight,keepaspectratio]{#1}%
  }{%
    \fcolorbox{openieCodeBorder}{openieCodeBg}{%
      \parbox{0.92\linewidth}{\footnotesize\ttfamily Figure asset embedded in archival PDF; interactive/vector source: \url{#3}}%
    }%
  }%
}

\title{\vspace{-1.2em}\textbf{Metabolic Intelligence: Budget-Native Compute from Tag to Campus}}
\author{\textbf{David Charlot, Open Interface Engineering}\\[0.15em]
{\small Product white paper $\cdot$ Open Interface Engineering $\cdot$ \href{https://research.openie.dev}{research.openie.dev}}}
\date{\small Revision \texttt{1561a2a} (22 November 2026) $\cdot$ Web: \href{https://research.openie.dev/products/mei/}{research.openie.dev/products/mei/} $\cdot$ Archival LaTeX: \href{https://research.openie.dev/pdfs/latex/mei.pdf}{/pdfs/latex/mei.pdf} (\href{https://research.openie.dev/pdfs/latex/mei.tex}{.tex})}

\begin{document}
\maketitle
\vspace{-0.6em}

% OpenIE Measurement Honesty Callout Box
\noindent\fcolorbox{openieAmberBorder}{openieAmberBg}{%
  \parbox{\dimexpr\linewidth-2\fboxsep-2\fboxrule\relax}{%
    \small\color{openieDark}%
    \textbf{Measurement.} Product white paper, kept apart from the research catalog; not a journal final. Combines emulator instruction counts, datasheet energy schedules (\texttt{est\_j}), a grid-following campus protocol test, and the open \textbf{\texttt{MEI-JCP-Tag-v1}} sovereign hardware reference package: a \texttt{DRC}-certified (\textbf{\texttt{0} violations}) \texttt{32 $\times$ 24 mm} 2-layer \texttt{STM32WBA52} + \texttt{TI INA238} + \texttt{FRAM} PCB (Gerber \texttt{RS-274X}, KiCad, BOM/CPL), a \texttt{36 $\times$ 28 $\times$ 7.5 mm} 3D NURBS B-Rep enclosure (\texttt{STEP}/\texttt{STL}), and a \textbf{\texttt{10,000}-wakeup} \texttt{CR2032} electro-thermal \texttt{JCP} \texttt{hold}/\texttt{settle} verification preserving \texttt{100\%} of critical alerts at \texttt{-20 \textdegree C} across \texttt{2,680} \texttt{Ed25519} canonical \texttt{det-CBOR} \texttt{JCR-1} receipts. Host build counters are \texttt{reported\_j}; physical tag bench power has not been measured (\texttt{board\_metered = false}, \texttt{measured\_j = None}). No fabricated citations.
  }%
}

\vspace{6pt}
% OpenIE Amber Left-Rule Summary Quote
\noindent{\color{openieAmber}\vrule width 3pt}\hspace{8pt}%
\parbox{\dimexpr\linewidth-11pt\relax}{%
  \small\itshape\color{openieSlate}%
  Computer intelligence designed to run inside a fixed energy budget, measured on a cold-chain sensor tag and modeled for a grid-following data center, with every number labeled by how it was obtained.
}
\vspace{6pt}

{
\small
\setcounter{tocdepth}{2}
\tableofcontents
}
\vspace{4pt}
\hrule
\vspace{8pt}

\section{Abstract}\label{abstract}

Metabolic Intelligence (MEI) is computer intelligence designed to run
inside a fixed energy budget. This paper defines it, reports emulator,
datasheet, and computer-aided design (CAD) verification results for a
cold-chain sensor tag, and reports a protocol-level test of a data
center that follows grid power caps. Every number carries the method
that produced it.

\textbf{In plain words}

\begin{enumerate}
\def\labelenumi{\arabic{enumi}.}
\tightlist
\item
  Give a computer a fixed energy budget, and design it to give the best
  answer it can inside that budget.
\item
  A tiny sensor tag that checks food freshness can then run for about
  eight years on a coin battery at room temperature, and under
  cold-chain voltage sag a signed budget gate refuses routine radio
  beacons so safety-critical alerts still get through.
\item
  A data center can cut its power use when the electricity grid asks, by
  pausing or shrinking work, without giving wrong answers.
\end{enumerate}

\section{1. Introduction}\label{introduction}

A computing system usually treats energy as a cost reported after the
work. Here energy is a condition stated before the work: a budget \(B\),
in joules (J) for a task or in watts (W) for a sustained load. The
design question is which answers can be obtained inside \(B\), and how
to obtain them at the least energy.

Three established results set the frame.

\begin{itemize}
\tightlist
\item
  \textbf{Landauer (1961).} Erasing one bit at temperature \(T\)
  dissipates at least \(k_B T \ln 2\), with
  \(k_B = 1.380649 \times 10^{-23}\) J/K. At \(T = 300\) K this is
  \(2.87 \times 10^{-21}\) J.
\item
  \textbf{Horowitz (2014).} Practical CMOS operations cost many orders
  of magnitude more: about 0.9 pJ for a 32-bit floating-point add and
  about 640 pJ for a 32-bit DRAM read at 45 nm. Data movement, not
  arithmetic, dominates.
\item
  \textbf{Howard (1966).} The value of information (VoI) of an
  observation is the expected gain in a decision's payoff from making it
  before deciding, minus its cost. An observation with VoI ≤ 0 is not
  worth making.
\end{itemize}

Together they say that energy per answer has a physical floor far below
current practice, that most of the gap is data movement, and that an
observation or computation is worth its joules only when it can change a
decision.

The paper studies two scales.

{\def\LTcaptype{none} % do not increment counter
{\footnotesize\setlength{\tabcolsep}{3.8pt}\renewcommand{\arraystretch}{1.15}
\begin{longtable}[]{@{}
  >{\raggedright\arraybackslash}p{(\linewidth - 6\tabcolsep) * \real{0.1905}}
  >{\raggedright\arraybackslash}p{(\linewidth - 6\tabcolsep) * \real{0.1905}}
  >{\raggedright\arraybackslash}p{(\linewidth - 6\tabcolsep) * \real{0.3095}}
  >{\raggedright\arraybackslash}p{(\linewidth - 6\tabcolsep) * \real{0.3095}}@{}}
\toprule\noalign{}
\begin{minipage}[b]{\linewidth}\raggedright
\textbf{Scale}
\end{minipage} & \begin{minipage}[b]{\linewidth}\raggedright
\textbf{Budget unit}
\end{minipage} & \begin{minipage}[b]{\linewidth}\raggedright
\textbf{Controls}
\end{minipage} & \begin{minipage}[b]{\linewidth}\raggedright
\textbf{Measured outcome}
\end{minipage} \\
\midrule\noalign{}
\endhead
\bottomrule\noalign{}
\endlastfoot
Sensor tag & µW sustained; nJ to mJ per event & Heater duty, MCU sleep,
classifier early exit, two-phase radio budget gate & Battery life;
accuracy under sequential learning; cold-chain alert survival \\
Data center (100 MW, grid-capped) & MW import cap & Pause, speed, model
choice, effort, battery, shed & Minutes over cap; share of interactive
requests served \\
\end{longtable}
}
}

\section{2. Definitions}\label{definitions}

{\def\LTcaptype{none} % do not increment counter
{\footnotesize\setlength{\tabcolsep}{3.8pt}\renewcommand{\arraystretch}{1.15}
\begin{longtable}[]{@{}
  >{\raggedright\arraybackslash}p{(\linewidth - 2\tabcolsep) * \real{0.2708}}
  >{\raggedright\arraybackslash}p{(\linewidth - 2\tabcolsep) * \real{0.7292}}@{}}
\toprule\noalign{}
\begin{minipage}[b]{\linewidth}\raggedright
\textbf{Term}
\end{minipage} & \begin{minipage}[b]{\linewidth}\raggedright
\textbf{Definition}
\end{minipage} \\
\midrule\noalign{}
\endhead
\bottomrule\noalign{}
\endlastfoot
Budget \(B\) & The energy (J) or power (W) available to a task or a
site, fixed before the work. \\
Utility \(U(a)\) & A stated score of answer \(a\): accuracy, a
completeness test, or a served request. \\
Completeness predicate \(C(z)\) & A written test that maps a task state
\(z\) to 0 or 1. The task is complete when \(C(z) = 1\). \\
Obtain mechanism \(m\) & How an input to a decision is acquired: Lookup
(read a stored value), Formula (evaluate a closed form), Solver (run a
search or optimizer), Model (run a learned model), Guess (use a
prior). \\
Binding classifier & A recognition unit with a binding site (a template
over input features), an abundance (a learned weight), a context switch
that gates which units are active, a Hill exponent for cooperative
response, and a counted cost per binding check. Tables use the short
label ``enzyme''. A bit-level variant uses no multiplications. \\
Metabolic schedule & A duty cycle for sensing and heating that meets a
monitoring objective at minimum average power. \\
Grid-following service & A control program that receives grid signals
and applies ordered controls so site import stays under the cap. \\
\end{longtable}
}
}

\textbf{Evidence classes.} Every number in this paper carries one of
these labels.

{\def\LTcaptype{none} % do not increment counter
{\footnotesize\setlength{\tabcolsep}{3.8pt}\renewcommand{\arraystretch}{1.15}
\begin{longtable}[]{@{}
  >{\raggedright\arraybackslash}p{(\linewidth - 4\tabcolsep) * \real{0.1919}}
  >{\raggedright\arraybackslash}p{(\linewidth - 4\tabcolsep) * \real{0.4646}}
  >{\raggedright\arraybackslash}p{(\linewidth - 4\tabcolsep) * \real{0.3434}}@{}}
\toprule\noalign{}
\begin{minipage}[b]{\linewidth}\raggedright
\textbf{Label}
\end{minipage} & \begin{minipage}[b]{\linewidth}\raggedright
\textbf{Meaning}
\end{minipage} & \begin{minipage}[b]{\linewidth}\raggedright
\textbf{What it can support}
\end{minipage} \\
\midrule\noalign{}
\endhead
\bottomrule\noalign{}
\endlastfoot
real protocol & Exchanged messages with a real server or broker &
Interoperability, not watts \\
emulator-measured & Instruction counts from the compiled binary run
under an emulator & Relative cost, not board joules \\
datasheet & Vendor electrical tables & Inputs to a schedule model \\
cad-verified & Programmatic electronic design automation, design-rule
check, solid geometry, and electro-thermal simulation & Fabrication
readiness and modeled derating, not bench \texttt{measured\_j} \\
host-counter (\texttt{reported\_j}) & Package energy register delta
(Linux Running Average Power Limit, RAPL) on the CAD build workstation &
Build-pipeline host compute reporting, not tag board joules \\
measured-kernel model & Throughput from measured GPU kernels combined
with published power & Energy-per-answer estimates with bands \\
looked-up & Published medians (for example served tokens per second) &
Baselines \\
calculation & Algebra from stated inputs & Consequences of the inputs \\
plant model & A simulator of site power and thermal behavior & Scenario
comparisons under stated assumptions \\
\end{longtable}
}
}

No figure in this paper is a reading from an external energy meter on a
fabricated tag or GPU rack. The page flag \texttt{board\_metered} is
false, and \texttt{measured\_j} remains absent (\texttt{null})
throughout.

\section{\texorpdfstring{3. The obtain cost
\(J(m \mid q)\)}{3. The obtain cost J(m \textbackslash mid q)}}\label{the-obtain-cost-jm-mid-q}

A decision \(y\) depends on an input quantity \(q\). Mechanism \(m\)
obtains \(q\) with energy \(E_m\) (J), a standard deviation \(\sigma_m\)
in relative units, a staleness rate \(\kappa_q\) (relative variance per
second) and an age \(t\) (s). Let
\(S_q = \partial \ln y / \partial \ln q\) be the elasticity of the
decision with respect to the input, and let \(\lambda\) (J) be the
energy at stake in the decision. To second order, the expected loss from
an input error of relative variance \(v\) is
\(\tfrac{1}{2}\lambda S_q^2 v\). The obtain cost is energy plus that
loss:

\[J(m \mid q) = E_m + \tfrac{1}{2}\,\lambda\, S_q^2\, \mathrm{var}_m(t), \qquad \mathrm{var}_m(t) = (\sigma_m^2 + \kappa_q t)\,(1 + 3 p_m),\]

and the chosen mechanism is \(m^* = \arg\min_m J(m \mid q)\). The factor
\((1 + 3p_m)\) inflates variance by the probability \(p_m\) that the
mechanism is applied outside its valid range; the coefficient 3 is a
modeling choice, stated as such. \(S_q\) is computed by central
differences on the decision function, not assumed.

\textbf{Lemma A (dominance).} If \(m_1\) and \(m_2\) obtain the same
\(q\) with \(E_{m_1} < E_{m_2}\) and
\(\mathrm{var}_{m_1} \le \mathrm{var}_{m_2}\), then
\(J(m_1 \mid q) < J(m_2 \mid q)\) for all \(\lambda, S_q \ge 0\).
\emph{Proof:} both terms are weakly smaller for \(m_1\) and the first
strictly. So a table lookup beats a calculation of the same value, and a
closed form beats a learned model of the same value.

\textbf{Lemma B (precision to buy).} Setting the loss term equal to a
tolerated excess cost \(\varepsilon\) (J) gives
\(\tfrac{1}{2}\lambda S_q^2 \sigma^2 = \varepsilon\), so
\(\sigma^* = \sqrt{2\varepsilon / (\lambda S_q^2)}\). Precision finer
than \(\sigma^*\) costs energy without reducing the loss below
\(\varepsilon\).

\textbf{Lemma C (refresh interval).} With staleness alone, the loss
grows as \(\tfrac{1}{2}\lambda S_q^2 \kappa_q t\) (J). Refreshing when
that loss equals the cost \(E\) of a new reading gives
\(t^* = 2E / (\lambda S_q^2 \kappa_q)\) (s). A sensor's sampling
interval and a cache's lifetime are the same quantity at different
\(\lambda\).

The estimator of \(J\) must itself be cheap: it is evaluated by Lookup
or Formula with a fixed energy cap, and when the cap is reached the
default mechanism is used.

\section{4. Sensor tag}\label{sensor-tag}

\subsection{4.1 Classifier cost}\label{classifier-cost}

Instruction counts per decision on an Arm Cortex-M4 with 12 input
features (emulator-measured):

{\def\LTcaptype{none} % do not increment counter
{\footnotesize\setlength{\tabcolsep}{3.8pt}\renewcommand{\arraystretch}{1.15}
\begin{longtable}[]{@{}
  >{\raggedright\arraybackslash}p{(\linewidth - 2\tabcolsep) * \real{0.6042}}
  >{\raggedright\arraybackslash}p{(\linewidth - 2\tabcolsep) * \real{0.3958}}@{}}
\toprule\noalign{}
\begin{minipage}[b]{\linewidth}\raggedright
\textbf{Classifier}
\end{minipage} & \begin{minipage}[b]{\linewidth}\raggedright
\textbf{Instructions per decision}
\end{minipage} \\
\midrule\noalign{}
\endhead
\bottomrule\noalign{}
\endlastfoot
Binding classifiers with family index and early exit & 1,723 \\
int8 multilayer perceptron (MLP) & 2,451 \\
Binding classifiers, first version without early exit & 6,948 \\
\end{longtable}
}
}

\textbf{Energy per decision (calculation).} The STM32U3 datasheet gives
15.5 µA/MHz at 3.3 V, so one clock cycle costs
\(15.5 \times 10^{-6}\,\text{A} \times 3.3\,\text{V} / 10^6\,\text{Hz} = 51\)
pJ. At about one instruction per cycle, 1,723 instructions cost about
\(1{,}800 \times 51\,\text{pJ} \approx 92\) nJ.

\subsection{4.2 Accuracy}\label{accuracy}

{\def\LTcaptype{none} % do not increment counter
{\footnotesize\setlength{\tabcolsep}{3.8pt}\renewcommand{\arraystretch}{1.15}
\begin{longtable}[]{@{}
  >{\raggedright\arraybackslash}p{(\linewidth - 4\tabcolsep) * \real{0.3394}}
  >{\raggedright\arraybackslash}p{(\linewidth - 4\tabcolsep) * \real{0.4312}}
  >{\raggedright\arraybackslash}p{(\linewidth - 4\tabcolsep) * \real{0.2294}}@{}}
\toprule\noalign{}
\begin{minipage}[b]{\linewidth}\raggedright
\textbf{Test}
\end{minipage} & \begin{minipage}[b]{\linewidth}\raggedright
\textbf{Result}
\end{minipage} & \begin{minipage}[b]{\linewidth}\raggedright
\textbf{Evidence}
\end{minipage} \\
\midrule\noalign{}
\endhead
\bottomrule\noalign{}
\endlastfoot
Same simulated features, 3 seeds & Binding classifiers 82.9 / 86.9 /
88.7\%; MLP 82.2 / 78.0 / 87.4\% & Real protocol on simulated
features \\
Late life (days 90 to 180) with drift-invariant features & LDA 94.9\%,
MLP 93.4\%, binding classifiers 90.0\% (80.8\% without these features),
Tsetlin machine 89.1\% & Simulation \\
SmellNet, 50 foods learned in 5 groups, no revisiting earlier groups &
Binding classifiers 77.7\%; MLP without replay or regularization 18.4\%
& Public dataset \\
SmellNet across days & MLP 85.5\%; binding classifiers 79.8\% (same on
Cortex-M4) & Public dataset \\
\end{longtable}
}
}

The 18.4\% MLP is the catastrophic-forgetting floor, not a
continual-learning baseline. The fair comparison runs elastic weight
consolidation (Kirkpatrick et al., 2017), experience replay with a small
buffer (Rolnick et al., 2019) and iCaRL (Rebuffi et al., 2017) under the
same memory and energy budget, scored for retained accuracy and energy
per inference under MLPerf Tiny rules (Banbury et al., 2021). That
comparison is specified and not yet run. On cross-day SmellNet the MLP
is more accurate. Drift-invariant features add 9 to 11 points of
late-life accuracy to binding classifiers, bit-level binding classifiers
and Tsetlin machines. A chicken-spoilage dataset (Anwar and Anwar, 2025)
was used internally; it has no licence for publication, so no result
from it appears here.

\subsection{4.3 Power budget and battery
life}\label{power-budget-and-battery-life}

One reading of a metal-oxide gas sensor costs 0.5 to 27 mJ depending on
the part (datasheets). The schedule takes about 9.6 readings per day.

{\def\LTcaptype{none} % do not increment counter
{\footnotesize\setlength{\tabcolsep}{3.8pt}\renewcommand{\arraystretch}{1.15}
\begin{longtable}[]{@{}
  >{\raggedright\arraybackslash}p{(\linewidth - 4\tabcolsep) * \real{0.4062}}
  >{\raggedright\arraybackslash}p{(\linewidth - 4\tabcolsep) * \real{0.2292}}
  >{\raggedright\arraybackslash}p{(\linewidth - 4\tabcolsep) * \real{0.3646}}@{}}
\toprule\noalign{}
\begin{minipage}[b]{\linewidth}\raggedright
\textbf{Term}
\end{minipage} & \begin{minipage}[b]{\linewidth}\raggedright
\textbf{Average power}
\end{minipage} & \begin{minipage}[b]{\linewidth}\raggedright
\textbf{Evidence}
\end{minipage} \\
\midrule\noalign{}
\endhead
\bottomrule\noalign{}
\endlastfoot
Heater (Bosch BME690) & 3.0 µW & Datasheet and schedule model \\
MCU sleep (ST STM32U3) & 6.9 µW & Datasheet \\
Total & 9.9 µW & Sum \\
Same tag with Nordic nRF54L15 sleep & 5.1 µW & Datasheet \\
\end{longtable}
}
}

\textbf{Check of the heater term.}
\(3.0\,\mu\text{W} \times 86{,}400\,\text{s/day} / 9.6\,\text{readings/day} = 27\)
mJ per reading, the upper end of the datasheet range. Classification at
92 nJ × 9.6 per day averages \(1.0 \times 10^{-11}\) W, negligible. At
this schedule sleep current, not heating or classification, is the
largest term, so the MCU is chosen for sleep current.

\textbf{Battery life (calculation).} A CR2032 cell holds about 235 mAh
to 2.0 V at 21 °C at 3.0 V nominal (Energizer datasheet), so
\(0.235\,\text{A·h} \times 3600\,\text{s/h} \times 3.0\,\text{V} = 2.54\)
kJ. The tag uses
\(9.9 \times 10^{-6}\,\text{W} \times 3.156 \times 10^7\,\text{s/yr} = 312\)
J per year, so the life is \(2540 / 312 = 8.1\) years without
self-discharge and about 7.5 years with 1\% per year self-discharge.
Lithium coin cells deliver less below 21 °C; Section 4.4 models that
electro-thermal sag across \(-20\text{ °C}\), \(+4\text{ °C}\), and
\(+25\text{ °C}\) and gates radio transmissions against it, while
physical bench discharge at cold-chain temperature remains a measurement
still to make. The 3.0 µW heater average implies a very low duty cycle
for a sensor that draws milliwatts when on. Whether sensing response
holds at that duty is also a measurement still to make.

\subsection{4.4 Reference hardware, cold-chain derating, and
cryptographic budget
receipts}\label{reference-hardware-cold-chain-derating-and-cryptographic-budget-receipts}

To verify that the tag architecture closes electrically, mechanically,
and under low-temperature coin-cell impedance rise, this study generates
an open reference package (\texttt{MEI-JCP-Tag-v1}) programmatically
through \texttt{openie-cad} (\texttt{cad-eda}, \texttt{cad-router},
\texttt{cad-kernel}, \texttt{cad-sim}, \texttt{cad-mfg},
\texttt{cad-interop}, and \texttt{cad-joules}) paired with the Joule
Context Protocol (\texttt{JCP}, draft specification \texttt{JCR-1}) and
the System-and-Task-Boundary Input/Output (\texttt{SATB-IO}) provenance
lattice (\texttt{jcp-\allowbreak{}satb-\allowbreak{}cad-\allowbreak{}spine}). The board is an exported
fabrication package, not yet physically assembled on a bench
(\texttt{board\_\allowbreak{}metered:\ false}, \texttt{measured\_\allowbreak{}j:\ null}).


\begin{center}
\openiefig{fig/mei-fig01-mei-jcp-tag-pcb.png}{0.94}{https://research.openie.dev/figures/mei/mei-jcp-tag-pcb.svg}
\end{center}


\textbf{2-layer printed circuit board (PCB) and 3D enclosure
(\texttt{cad-verified}).} The \texttt{MEI-JCP-Tag-v1} board measures
\(32.0\text{ mm} \times 24.0\text{ mm}\) on 2-layer 1.6 mm FR-4
(\texttt{0.18\ mm} signal traces, \texttt{0.\allowbreak{}30\ mm\ /\allowbreak{}\ 0.\allowbreak{}56\ mm} plated
through-vias). Ten components realize the sensing, inline rail
accounting, and hash-chained receipt storage loop:

\begin{enumerate}
\def\labelenumi{\arabic{enumi}.}
\tightlist
\item
  \texttt{U1} (\texttt{STM32WBA52CGU}): ultra-low-power microcontroller
  and long-range Bluetooth Low Energy radio gate on \texttt{F.Cu}.
\item
  \texttt{U2} (\texttt{INA238AIDR}) + \texttt{R\_SHUNT}
  (\texttt{100\ mΩ}, \texttt{0.1\%}, \texttt{25\ ppm/°C}): 16-bit
  inter-integrated circuit (\texttt{I2C}) shunt monitor across the raw
  \texttt{CR2032} feed (\texttt{VBAT\_RAW} to \texttt{VDD\_SYS}) with
  hardware alert pin \texttt{JCP\_ALERT}.
\item
  \texttt{U3} (\texttt{MB85RS256B}): 256-kilobit serial peripheral
  interface (\texttt{SPI}) ferroelectric random-access memory
  (\texttt{FRAM}) storing the hash-chained \texttt{JCR-1} receipt ring
  across brownouts.
\item
  \texttt{RT1} + \texttt{R\_BIAS} (\texttt{10\ kΩ} negative temperature
  coefficient, \texttt{NTC}, thermistor divider), \texttt{C1}
  (\texttt{100\ nF}), \texttt{C2} (\texttt{1\ µF}), \texttt{J1} (5-pin
  serial wire debug and receipt-dump universal asynchronous
  receiver-transmitter header), and \texttt{BT1}
  (\texttt{Keystone\ 3034} \texttt{CR2032} holder on \texttt{B.Cu}).
\end{enumerate}

\texttt{cad-router} decomposes all 14 nets into 32 Steiner connections
and routes 100.0\% (\texttt{135} copper segments totaling
\texttt{280.23\ mm}, \texttt{14} vias). \texttt{cad-eda} Design Rule
Check (\texttt{DRC}) and certification pass with \textbf{0 errors and 0
warnings (\texttt{CERTIFIED})}, round-trip verified through
\texttt{.kicad\_pcb} import/export alongside RS-274X Gerbers, Excellon
\texttt{.drl}, bill of materials (\texttt{bom.csv}), and component
placement (\texttt{cpl.csv}). In \texttt{cad-kernel}, the snap-fit
vented enclosure (\texttt{36.\allowbreak{}0\ mm\ ×\ 28.\allowbreak{}0\ mm\ ×\ 7.\allowbreak{}5\ mm} main shell,
\texttt{41.5\ mm} overall X span with a \texttt{Ø4.0\ mm} lanyard
mounting ear, convective airflow slot over \texttt{RT1}, and four
internal standoffs) evaluates as a closed outward-facing Boundary
Representation (\texttt{B-Rep}) solid of volume \texttt{2,702.12\ mm³}
(\texttt{2.837\ g} in acrylonitrile butadiene styrene, \texttt{ABS}),
exported to ISO-10303 \texttt{STEP} (\texttt{174,694\ B}) and binary
\texttt{STL} (\texttt{852} triangles).

\textbf{Cold-chain \texttt{CR2032} electrochemistry and
\texttt{10,000}-wakeup \texttt{JCP} two-phase gate
(\texttt{cad-verified}, \texttt{est\_j}).} Evaluating
\texttt{cad-\allowbreak{}sim::battery} on a calibrated \texttt{CR2032} model
(\texttt{3.0\ V}, \texttt{225\ mAh} = \texttt{2,430\ J} nominal,
\(R_{\text{int}}(50\%\text{ SOC}, 25\text{ °C}) = 12.0\ \Omega\),
Peukert exponent \(1.12\)) quantifies why cold-chain tags fail when
routine radio bursts are unbudgeted:

{\def\LTcaptype{none} % do not increment counter
{\scriptsize\setlength{\tabcolsep}{3.0pt}\renewcommand{\arraystretch}{1.14}
\begin{longtable}[]{@{}
  >{\raggedright\arraybackslash}p{(\linewidth - 12\tabcolsep) * \real{0.0843}}
  >{\raggedright\arraybackslash}p{(\linewidth - 12\tabcolsep) * \real{0.1461}}
  >{\raggedright\arraybackslash}p{(\linewidth - 12\tabcolsep) * \real{0.1517}}
  >{\raggedright\arraybackslash}p{(\linewidth - 12\tabcolsep) * \real{0.1798}}
  >{\raggedright\arraybackslash}p{(\linewidth - 12\tabcolsep) * \real{0.1685}}
  >{\raggedright\arraybackslash}p{(\linewidth - 12\tabcolsep) * \real{0.1573}}
  >{\raggedright\arraybackslash}p{(\linewidth - 12\tabcolsep) * \real{0.1124}}@{}}
\toprule\noalign{}
\begin{minipage}[b]{\linewidth}\raggedright
\textbf{Ambient \(T\)}
\end{minipage} & \begin{minipage}[b]{\linewidth}\raggedright
\textbf{\(R_{\text{int}}\) at 50\% SOC (\texttt{cad-sim})}
\end{minipage} & \begin{minipage}[b]{\linewidth}\raggedright
\textbf{12 mA TX terminal voltage (50\% SOC)}
\end{minipage} & \begin{minipage}[b]{\linewidth}\raggedright
\textbf{Usable energy to 2.00 V at 12 mA (\texttt{est\_j})}
\end{minipage} & \begin{minipage}[b]{\linewidth}\raggedright
\textbf{\texttt{10,000}-wakeup routine beacons
(\texttt{allow\ /\allowbreak{}\ refused})}
\end{minipage} & \begin{minipage}[b]{\linewidth}\raggedright
\textbf{Safety-critical alerts (\texttt{allow\ /\allowbreak{}\ attempts})}
\end{minipage} & \begin{minipage}[b]{\linewidth}\raggedright
\textbf{Verified \texttt{JCR-1} receipts}
\end{minipage} \\
\midrule\noalign{}
\endhead
\bottomrule\noalign{}
\endlastfoot
\(+25\text{ °C}\) & \(12.00\ \Omega\) & \(2.669\text{ V}\)
(\(144.0\text{ mV}\) sag) & \(2{,}193.12\text{ J}\) (\(90.25\%\)) &
\(500 / 0\) & \(40 / 40\) (\(100.0\%\)) & \(1{,}080\) \\
\(+4\text{ °C}\) & \(25.72\ \Omega\) & \(2.504\text{ V}\)
(\(308.7\text{ mV}\) sag) & \(2{,}059.23\text{ J}\) (\(84.74\%\)) &
\(430 / 70\) & \(40 / 40\) (\(100.0\%\)) & \(1{,}010\) \\
\(-20\text{ °C}\) & \(71.78\ \Omega\) & \(1.951\text{ V}\)
(\(861.3\text{ mV}\) sag) & \(0.00\text{ J}\) (unbuffered pulse
\(< 2.00\text{ V}\)) & \(10 / 490\) & \(40 / 40\) (\(100.0\%\)) &
\(590\) \\
\end{longtable}
}
}

At \(-20\text{ °C}\), Arrhenius impedance growth raises
\(R_{\text{int}}\) nearly six-fold (\(71.78\ \Omega\) at 50\% state of
charge), pulling an unbuffered \(12\text{ mA}\) radio burst down to
\(1.951\text{ V}\) (below the \(2.00\text{ V}\) brownout threshold).
Sharding the mission grant via \texttt{jcp\_\allowbreak{}grant::shard} into an
attenuated routine-beacon child grant and an isolated \(35\text{ mJ}\)
safety-excursion child grant --- enforced before every transmission via
\texttt{jcp\_\allowbreak{}runtime::Runtime::hold\_\allowbreak{}chained} and resolved afterward via
\texttt{settle\_chained} --- deterministically refuses \(490\) of
\(500\) non-essential routine beacons under \(-20\text{ °C}\) derating
while preserving \textbf{40 of 40 (100.0\%)} safety-critical thermal
excursion alerts across all \(10{,}000\) wakeups. Every run passes
Ed25519 \texttt{COSE\_Sign1} verification
(\texttt{jcp\_\allowbreak{}receipt::verify\_\allowbreak{}log}) and crash-recovery ledger
reconstruction (\texttt{jcp\_\allowbreak{}receipt::replay\_\allowbreak{}ledger}).


\begin{center}
\openiefig{fig/mei-fig02-jcp-satb-cad-spine.png}{0.94}{https://research.openie.dev/figures/mei/jcp-satb-cad-spine.svg}
\end{center}


\textbf{Provenance meet-semilattice and build-pipeline accounting
(\texttt{reported\_j} vs \texttt{est\_j}).} Across the 7-stage
\texttt{openie-cad} build on a 32-core Linux workstation
(\texttt{6.93\ s} wall time), \texttt{jcp-\allowbreak{}satb-\allowbreak{}cad-\allowbreak{}spine} enforces a
monotone meet-semilattice over \texttt{280\ ×\ 280\ =\ 78,\allowbreak{}400} tested
\texttt{(\allowbreak{}SATB\ boundary\ ×\ attribution\ ×\ mechanism)} pairs
(\texttt{0} silent provenance or rail-coverage promotions),
\texttt{SATB} Falsifier 6 (\texttt{\textgreater{}\ 5.00\%} pre/post idle
baseline drift omits \texttt{satb\_\allowbreak{}derived\_\allowbreak{}marginal\_\allowbreak{}uj} from the
signed canonical deterministic Concise Binary Object Representation,
\texttt{det-CBOR}, core while retaining
\texttt{satb\_\allowbreak{}raw\_\allowbreak{}counter\_\allowbreak{}uj}), and \texttt{SATB} Step 7
(\texttt{\textbar{}sum(phases)\ -\ window\textbar{}\ \textless{}=\ 1\ µJ},
observed residual \texttt{0\ µJ}). Host Linux \texttt{RAPL}
(\texttt{package-0}) recorded \texttt{380.20\ J} (\texttt{reported\_j},
vendor-opaque \texttt{HwShunt} \texttt{Partial}) alongside
\texttt{242.07\ J} analytical \texttt{est\_j} (\texttt{45.15\ J}
embodied FR-4 + ABS material estimate, \texttt{14.50\ J} surface-mount
assembly batch estimate, and \texttt{182.42\ J} analytical compute
estimate), sealed in \texttt{8} verified \texttt{JCR-1} receipts
(\texttt{head\ =\ d5b3827b.\allowbreak{}.\allowbreak{}.\allowbreak{}01b2cd}), while physical tag
\texttt{measured\_j} stays \texttt{null}.

\section{5. Data center energy per
answer}\label{data-center-energy-per-answer}

\textbf{Method.} Energy per answer (Wh) = output tokens × power per GPU
(kW) ÷ output tokens per second per GPU ÷ 3600. Throughput comes from
NVIDIA AI Configurator 0.12.0 (measured GPU kernels, SGLang, NVFP4).
Power per GPU is published throughput per chip divided by published
throughput per MW (InferenceX): GB300 2.12 kW, B300 1.90 kW, B200 1.71
kW, Vera Rubin 3.30 kW, MI355X 2.09 kW. Served speeds per user are
Artificial Analysis medians: DeepSeek V4 Pro 79 tokens/s, V4 Flash 125,
Qwen3-VL 48.5.

{\def\LTcaptype{none} % do not increment counter
{\footnotesize\setlength{\tabcolsep}{3.8pt}\renewcommand{\arraystretch}{1.15}
\begin{longtable}[]{@{}
  >{\raggedright\arraybackslash}p{(\linewidth - 4\tabcolsep) * \real{0.3750}}
  >{\raggedright\arraybackslash}p{(\linewidth - 4\tabcolsep) * \real{0.3750}}
  >{\raggedright\arraybackslash}p{(\linewidth - 4\tabcolsep) * \real{0.2500}}@{}}
\toprule\noalign{}
\begin{minipage}[b]{\linewidth}\raggedright
\textbf{Workload (GB300)}
\end{minipage} & \begin{minipage}[b]{\linewidth}\raggedright
\textbf{Energy per answer}
\end{minipage} & \begin{minipage}[b]{\linewidth}\raggedright
\textbf{Evidence}
\end{minipage} \\
\midrule\noalign{}
\endhead
\bottomrule\noalign{}
\endlastfoot
Chat, DeepSeek V4 Pro (2k tokens in, 500 out) & 0.32 Wh at 50
tokens/s/user; 0.92 at 79; 4.57 at 159 & Measured-kernel model \\
Same, fused estimate & 0.20 Wh (10th to 90th percentile 0.04 to 0.68) &
Model and looked-up \\
Chat, V4 Flash & 49 mWh at 50 tokens/s/user; 180 mWh at 159 &
Measured-kernel model \\
Agent turn, V4 Flash (32k in, 28k cached, 1k out) & 0.16 Wh at 96
tokens/s/user & Measured-kernel model \\
Video clip, 8 s at 720p (Wan 2.2) & 390 Wh & Model \\
\end{longtable}
}
}

Energy per answer rises steeply with speed per user because higher
per-user speed lowers the batch each GPU serves. Matching stream speed
to human reading speed raises work per MW by an estimated 1.92× on GB300
(band 1.19 to 2.93) and 2.0× on Vera Rubin (band 1.2 to 3.4). Routing
most chat to the smaller model plus speed matching gives an estimated
2.3 to 4.4×. Speculative decoding with a modeled accept length of 2.8
lowers energy per output token by about 1.55×, after which speed
matching adds only 1.4 to 1.8×. These are estimates with bands; the
measurement that would narrow them most is a day of metered energy on
one GPU pool.

\section{6. Grid-following data
center}\label{grid-following-data-center}

\subsection{6.1 Controls and their
order}\label{controls-and-their-order}

When the grid asks for a cut, deferrable work is acted on in a fixed
order: pause deferrable work, lower per-user speed, route to a smaller
model, lower reasoning effort, discharge battery, shed load. Interactive
and certified work runs to completion or is refused; it is not degraded.
A contract declares invariants (for example
\texttt{never\_\allowbreak{}lower\_\allowbreak{}precision}, \texttt{max\_deferral}), and shedding
is used only outside them. Routing to a smaller model changes which
model answers deferrable work; it is not assumed to be free of quality
loss (V4 Flash trails V4 Pro by about 1.5 points on GPQA).

\subsection{6.2 System}\label{system}

\begin{itemize}
\tightlist
\item
  \textbf{Grid input:} an OpenADR 3.1 client (OAuth2, polling, demand
  reports) tested against the open-source openleadr-rs server. Real
  protocol.
\item
  \textbf{Backstop:} a DSX Flex client over MQTT with CloudEvents and
  JWS ES256 signatures, tested against a real broker and schemas. Real
  protocol.
\item
  \textbf{Controls:} GPU power budgets in the Dynamo planner, a gateway
  that caps in-flight requests and sets routing, effort and speed tiers,
  job pause, Redfish node sleep, battery converter.
\item
  \textbf{Learning:} per-pool speed curves by recursive least squares; a
  bias correction on the power model (it learned +5 MW on a hot day).
\item
  \textbf{Plant:} a model of a 100 MW GB300 site. Plant model.
\end{itemize}

\subsection{6.3 Results}\label{results}

{\def\LTcaptype{none} % do not increment counter
{\footnotesize\setlength{\tabcolsep}{3.8pt}\renewcommand{\arraystretch}{1.15}
\begin{longtable}[]{@{}
  >{\raggedright\arraybackslash}p{(\linewidth - 6\tabcolsep) * \real{0.3299}}
  >{\raggedright\arraybackslash}p{(\linewidth - 6\tabcolsep) * \real{0.1753}}
  >{\raggedright\arraybackslash}p{(\linewidth - 6\tabcolsep) * \real{0.1856}}
  >{\raggedright\arraybackslash}p{(\linewidth - 6\tabcolsep) * \real{0.3093}}@{}}
\toprule\noalign{}
\begin{minipage}[b]{\linewidth}\raggedright
\textbf{Scenario}
\end{minipage} & \begin{minipage}[b]{\linewidth}\raggedright
\textbf{Minutes over cap}
\end{minipage} & \begin{minipage}[b]{\linewidth}\raggedright
\textbf{Interactive served}
\end{minipage} & \begin{minipage}[b]{\linewidth}\raggedright
\textbf{Note}
\end{minipage} \\
\midrule\noalign{}
\endhead
\bottomrule\noalign{}
\endlastfoot
A to C: 60/70\% cut, with notice, feedforward, or no notice & 0 & 100\%
& Notice mainly saves battery \\
D: 90/95\% cut, fast site (159 tokens/s/user) & 0 & 100\% & All
knowledge chat routed; effort low \\
E: 90/95\% cut without node sleep & 96 & 1\% & Idle GPUs alone exceed
the 10 MW cap \\
L to M: looked-up speeds, 60/70\% and 90/95\% cuts & 0 & 100\% / 52.6\%
& Speed already low cannot be lowered further \\
N: no notice, no battery, 10 min response & 7.5 & 99.8\% & Under cap
within 225 s \\
O: 60\% cut for 10 h & 0 & 100\% & Deadlines kept \\
P: new cap every 5 min for 3 h & 0 & 98.8\% & \\
Q: recovery ramp limited to 2 MW/min & 0 & 100\% & Maximum rise 2.9
MW/min against 29 to 47 unlimited \\
R: 90\% cut for 10 h & 0 & 13\% & Video deadlines take priority by
policy \\
\end{longtable}
}
}

In every run no control call failed, the server accepted every OpenADR
report, and no DSX message was rejected. With node sleep available, time
over the cap was zero in every scenario. The cost of holding the cap is
measured as availability: the interactive share served falls to 52.6\%,
13\% and 1\% in scenarios M, R and E. Refused and deferred work is
counted, not hidden.

\textbf{Lemma D (sleep is required at deep cuts).} If idle import
\(P_{\text{idle}}\) exceeds the cap \(P_{\text{cap}}\), no combination
of speed, routing or effort controls that leaves nodes powered meets the
cap, because those controls act only on power above idle. Scenario E is
the instance: about 10 MW cap, idle import above it, 96 minutes over.

\subsection{6.4 The stop rule}\label{the-stop-rule}

When \(C(z) = 1\) and further work has VoI ≤ 0, further energy buys
nothing on the stated objective. The system enacts that stop at four
points: do not run a model when Lookup or Formula closes the input; do
not stream faster than a human reads when the consumer is a human; pause
or shed when the import cap binds, in order of least service lost per
MW; skip a heater pulse when the expected change in the reading is below
threshold and the battery budget binds.

\section{7. Related work}\label{related-work}

\textbf{Energy-based and associative models.} Hopfield (1982) cast
associative memory as descent on an energy. Krotov and Hopfield (2016)
raised capacity with higher-order energies (Dense Associative Memory).
Ramsauer et al.~(2020) showed softmax attention is a modern Hopfield
update. Hoover et al.~(2023) built attention layers that minimize a
token energy, and Gladstone et al.~(2025) trained transformers that
predict by minimizing a learned energy. Bodnia and Hanin (2026) describe
a reasoning model scored by an energy over traces. In all of these the
``energy'' is a scalar objective, not joules.

\textbf{Latent world models.} LeCun (2022) proposed joint-embedding
predictive architectures. V-JEPA 2-AC (Assran et al., 2025) plans a
robot arm from image goals after large video pretraining plus under 62
hours of robot video. The no-pretraining property of the binding
classifiers in Section 4 applies only to that path; latent world models
depend on large pretraining.

\textbf{Efficient sequence models.} Mamba (Gu and Dao, 2023) and Mamba-2
(Dao and Gu, 2024) scale linearly in sequence length with constant-size
state.

\textbf{Edge classifiers.} Tsetlin machines (Granmo, 2018) learn
propositional clauses; an all-digital 65 nm accelerator reaches 8.6 nJ
per MNIST frame (Tunheim et al., 2025). Hyperdimensional computing
(Kanerva, 2009) computes on high-dimensional random vectors.

\textbf{Substrates.} Intel's Hala Point has 1,152 Loihi 2 processors and
a stated maximum of 2,600 W (Intel, 2024). Chen et al.~(2025) ran ALBERT
on a 14 nm analog in-memory chip. Ahmed et al.~(2025) ran ResNet, BERT
and an Atari agent on a photonic processor. No published neuromorphic
gas-sensing energy figure was found.

\textbf{Energy-aware serving.} PEARL, OmniRouter (2025) and energy-aware
reasoning-model routing (2026) choose which language model answers.
NVIDIA Dynamo 1.5 adds power-aware scaling, and DSX Flex (Emerald AI,
2026) offers site flexibility to the grid. The obtain cost of Section 3
differs in choosing among non-model mechanisms before any model runs.

\textbf{Price and energy.} The money price of fixed inference capability
has fallen sharply (Epoch AI, 2025). A falling price does not lower the
physical energy per operation, which Landauer and Horowitz bound from
below.

\section{8. Limits}\label{limits}

\begin{itemize}
\tightlist
\item
  No figure is from an external energy meter on a physically assembled
  tag or GPU pool (\texttt{board\_\allowbreak{}metered:\ false},
  \texttt{measured\_\allowbreak{}j:\ null}). Tag steady-state power is datasheet plus
  schedule model; instruction counts are emulator-measured; the
  \texttt{MEI-JCP-Tag-v1} 2-layer PCB, 3D enclosure, and
  \texttt{-\allowbreak{}20\ °C\ /\allowbreak{}\ +4\ °C\ /\allowbreak{}\ +25\ °C} coin-cell impedance curves are
  CAD-verified fabrication exports and electrochemical simulations
  (\texttt{est\_j}), with Linux \texttt{RAPL} host build counters kept
  strictly in \texttt{reported\_j}; data center power is a plant model
  driven through real protocols.
\item
  While Section 4.4 models \texttt{CR2032} Arrhenius impedance sag and
  verifies \texttt{10,000}-wakeup \texttt{JCP} radio refusal against it
  in simulation, physical chamber discharge of assembled
  \texttt{MEI-JCP-Tag-v1} boards and gas-sensor response at the stated
  heater duty remain bench measurements still to make.
\item
  The sequential-learning comparison uses an MLP without
  continual-learning methods; the fair comparison is specified and not
  run.
\item
  Energy-per-answer figures are estimates with 10th to 90th percentile
  bands.
\item
  No claim is made about water saved or about equivalence with an
  uncapped site serving the same workload.
\end{itemize}

\section{9. Conclusion}\label{conclusion}

Stating energy as a budget before the work turns several design choices
into calculations: which mechanism obtains an input (Section 3), how
often a sensor samples (Lemma C), which radio transmissions a cold-chain
tag admits when low-temperature cell impedance rises (Section 4.4), and
which controls a site applies under a grid cap (Section 6). On the tag
the room-temperature schedule gives about 8 years on a coin cell at 9.9
µW, and the DRC-clean \texttt{MEI-JCP-Tag-v1} reference package
(\texttt{32\ mm\ ×\ 24\ mm} 2-layer PCB + vented
\texttt{36\ mm\ ×\ 28\ mm\ ×\ 7.\allowbreak{}5\ mm} enclosure + Ed25519
\texttt{JCR-1} receipt ring) verifies in simulation that two-phase
\texttt{hold} / \texttt{settle} sharding preserves 100\% of
safety-critical excursion alerts when \(-20\text{ °C}\) voltage sag
forces routine beacons to be refused. At the site, controls applied in a
fixed order kept import under the cap in every modeled scenario with
node sleep, at a measured cost in interactive service. The next
measurements need an external energy meter: assembled
\texttt{MEI-JCP-Tag-v1} board shunts in a cold chamber, and a day of
energy on one GPU pool.

\section{Appendix A. Notation}\label{appendix-a.-notation}

{\def\LTcaptype{none} % do not increment counter
{\footnotesize\setlength{\tabcolsep}{3.8pt}\renewcommand{\arraystretch}{1.15}
\begin{longtable}[]{@{}
  >{\raggedright\arraybackslash}p{(\linewidth - 4\tabcolsep) * \real{0.2268}}
  >{\raggedright\arraybackslash}p{(\linewidth - 4\tabcolsep) * \real{0.5464}}
  >{\raggedright\arraybackslash}p{(\linewidth - 4\tabcolsep) * \real{0.2268}}@{}}
\toprule\noalign{}
\begin{minipage}[b]{\linewidth}\raggedright
\textbf{Symbol}
\end{minipage} & \begin{minipage}[b]{\linewidth}\raggedright
\textbf{Meaning}
\end{minipage} & \begin{minipage}[b]{\linewidth}\raggedright
\textbf{Unit}
\end{minipage} \\
\midrule\noalign{}
\endhead
\bottomrule\noalign{}
\endlastfoot
\(B\) & Energy or power budget & J or W \\
\(C(z)\) & Completeness predicate & 0 or 1 \\
\(q\), \(y\) & Input quantity and decision & \\
\(m\) & Obtain mechanism & \\
\(E_m\) & Energy of mechanism \(m\) & J \\
\(S_q\) & Elasticity \(\partial \ln y / \partial \ln q\) &
dimensionless \\
\(\lambda\) & Energy at stake & J \\
\(\sigma_m\) & Standard deviation of the obtained value & relative \\
\(\kappa_q\) & Staleness rate & relative variance per s \\
\(t\) & Age of the value & s \\
\(p_m\) & Probability of use outside the valid range & \\
\(J(m \mid q)\) & Obtain cost & J \\
\(P_{\text{cap}}\), \(P_{\text{idle}}\) & Grid cap and idle import &
W \\
\end{longtable}
}
}

\section{References}\label{references}

Ahmed, S. R., et al.~(2025). Universal photonic artificial intelligence
acceleration. \emph{Nature}.
\url{https://doi.org/10.1038/s41586-025-08854-x}

Assran, M., et al.~(2025). V-JEPA 2: Self-Supervised Video Models Enable
Understanding, Prediction and Planning. arXiv:2506.09985.

Banbury, C., et al.~(2021). MLPerf Tiny Benchmark. arXiv:2106.07597.

Bodnia, E., and Hanin, B. (2026). Energy-Based Models for Reasoning,
LLMs for the Interface.
\url{https://logicalintelligence.com/blog/energy-based-models-for-reasoning}

Bosch Sensortec BME690; Renesas ZMOD4410; Sensirion SGP41; ST STM32U3;
ST STM32WBA52; Nordic nRF54L15; Texas Instruments INA238; Fujitsu
MB85RS256B. Product datasheets, 2026.

Chen, A., et al.~(2025). Demonstration of transformer-based ALBERT model
on a 14nm analog AI inference chip. \emph{Nature Communications}.
\url{https://doi.org/10.1038/s41467-025-63794-4}

Dao, T., and Gu, A. (2024). Transformers are SSMs: Generalized Models
and Efficient Algorithms Through Structured State Space Duality.
arXiv:2405.21060.

Emerald AI (2026). DSX Flex deployment with NVIDIA and Silicon Valley
Power. \url{https://www.emeraldai.co/blog/nvidia-dsx-pilot-framework}

Energizer. CR2032 product datasheet.
\url{https://data.energizer.com/pdfs/cr2032.pdf}

Epoch AI (2025). LLM inference prices.
\url{https://epoch.ai/data-insights}

Gladstone, A., et al.~(2025). Energy-Based Transformers are Scalable
Learners and Thinkers. arXiv:2507.02092.

Granmo, O.-C. (2018). The Tsetlin Machine. arXiv:1804.01508.

Gu, A., and Dao, T. (2023). Mamba: Linear-Time Sequence Modeling with
Selective State Spaces. arXiv:2312.00752.

Hoover, B., et al.~(2023). Energy Transformer. NeurIPS.
arXiv:2302.07253.

Hopfield, J. J. (1982). Neural networks and physical systems with
emergent collective computational abilities. \emph{PNAS}
79(8):2554-2558.

Horowitz, M. (2014). Computing's energy problem (and what we can do
about it). ISSCC, 10-14.

Howard, R. A. (1966). Information Value Theory. \emph{IEEE Transactions
on Systems Science and Cybernetics} 2(1):22-26.

InferenceX (2026). Vera Rubin NVL72 agentic inference.
\url{https://inferencex.semianalysis.com/blog/vera-rubin-nvl72-agentic-inference}

Intel (2024). Hala Point neuromorphic system.
\url{https://newsroom.intel.com/artificial-intelligence/intel-builds-worlds-largest-neuromorphic-system-to-enable-more-sustainable-ai}

Kanerva, P. (2009). Hyperdimensional Computing. \emph{Cognitive
Computation} 1:139-159.

Kirkpatrick, J., et al.~(2017). Overcoming catastrophic forgetting in
neural networks. \emph{PNAS} 114(13):3521-3526.

Krotov, D., and Hopfield, J. J. (2016). Dense Associative Memory for
Pattern Recognition. arXiv:1606.01164.

Landauer, R. (1961). Irreversibility and Heat Generation in the
Computing Process. \emph{IBM Journal of Research and Development}
5(3):183-191.

LeCun, Y. (2022). A Path Towards Autonomous Machine Intelligence,
version 0.9.2. \url{https://openreview.net/pdf?id=BZ5a1r-kVsf}

NVIDIA (2026). Dynamo 1.5 documentation.
\url{https://docs.nvidia.com/dynamo}

OmniRouter (2025). arXiv:2502.20576.

OpenADR Alliance. OpenADR 3.1 specification; openleadr-rs reference
server.

Ramsauer, H., et al.~(2020). Hopfield Networks is All You Need.
arXiv:2008.02217.

Rebuffi, S.-A., et al.~(2017). iCaRL: Incremental Classifier and
Representation Learning. CVPR.

Rolnick, D., et al.~(2019). Experience Replay for Continual Learning.
NeurIPS. arXiv:1811.11682.

Transaction Science (2026). Joule Context Protocol (\texttt{JCP} /
\texttt{JCR-1}) draft specification.
\url{https://jcp.transaction.science/}

Tunheim, S. A., et al.~(2025). An All-digital 65-nm Tsetlin Machine
Image Classification Accelerator with 8.6 nJ per MNIST Frame.
arXiv:2501.19347.


\end{document}
